\documentclass[a4paper, 12pt]{article} 
%%%% choose one of the following 4 lines to set colour and/or  date style, 
%\usepackage{pa} % Day, num Month year, and colour (use lp2)
\usepackage{fancyhdr}
\usepackage[left=2.5cm,top=3cm,right=2.5cm,bottom=3cm,bindingoffset=0.5cm]{geometry}

\fancyhead{}
\fancyhead[RO,RE]{\bf Bruce A. Desmarais}
\fancyhead[LO,LE]{\bf Contributions to Diversity}


\begin{document} 
\pagestyle{fancy}

\vspace*{-.8cm}

My commitment to advancing diversity and inclusiveness in the academy was shaped by early-career experiences as a first-generation college graduate from a working-class background. Several incidents in graduate school and in my early years as an assistant professor left me with a class-based sense of exclusion. However, acknowledging my privilege as a white man who is a native English speaker, this early career sense of exclusion provided me with a deep concern for the challenges that those from underrepresented gender, racial and ethnic groups must face in their career pursuits. In my research, teaching, and service, I have striven to treat the advancement of inclusion as a foundational guiding value.

The most consistent and direct way that I have worked to advance diversity is through the use of positions of power and authority to assure that the spaces over which I have influence are diverse and inclusive. I will highlight a few examples of what this has looked like. For the last seven years I have been the coordinator of a speaker series---first in the political science department at Penn State, and now in the Center for Social Data Analytics. As I plan speaker invites each year, I assure that the majority of the speakers in the series are from underrepresented gender and/or racial and ethnic groups. When I served as president of the American Political Science Association's Political Networks Section, and later as host of the Political Networks Conference, I made sure that a majority of those individuals who were nominated to serve in section service roles, and the 18-person conference program committee, respectively, were from underrepresented gender and/or racial and ethnic groups. For the past four years I have served as director of graduate studies for Penn State's dual-title PhD program on Social Data Analytics. In this role, I am responsible for recruitment and chairing admissions. I have made diversity and inclusion a priority in my outreach efforts, and our dual-title program is the only data science graduate program that we have been able to identify in which the majority of students are women. 

In addition to using general leadership roles to advance diversity, I have sought out service roles that are specifically focused on the advancement of diversity. For the past two years I have served on the Justice, Equity, Diversity, and Inclusion (JEDI) committee of the Institute for Computational and Data Sciences (ICDS) at Penn State. Via this committee I have helped to organize events and allocate resources to enhance the use of data science in research JEDI topics, and in making data science at Penn State more inclusive. I served on the organizing committee for a Penn State Symposium on Harnessing the Data Revolution to Enhance Diversity. As the organizer of the event series at the Center for Social Data Analytics, I dedicated Center resources to co-sponsor this Symposium. In Summer 2020 I co-led an ad-hoc effort to raise a dialogue regarding race and graduate education at my PhD alma mater, UNC Chapel Hill. In Spring 2020 one of my fellow PhD alumni, who is Black, went public with several concerns regarding unequal treatment at UNC based on race. Soon after, several other fellow alumni from underrepresented racial backgrounds raised related concerns. I worked with other alumni to write a letter to the UNC Political Science Department requesting that we hold an event where current faculty and students could discuss these issues openly with PhD alumni. This event happened in Fall 2020. As a show of good faith, we also raised a substantial sum to donate to the UNC Graduate School's program on Diversity and Student Success. It is important that specific institutional efforts are made to advance the diversity and inclusiveness of the spaces in which we work, and equally important that this service does not disproportionately fall on those from underrepresented backgrounds.

Beyond my service work, I have striven to advance diversity through my practice of teaching and research. I teach four courses at Penn State (one undergraduate, and three graduate)---all are research methods courses. On each syllabus a majority of the authors listed in the assigned readings are from underrepresented gender, racial and ethnic groups. For at least two weeks in each course, the application readings covered address the areas of gender and politics and race and ethnic politics. I conduct team-based research, involving a variety of faculty, postdoctoral scholars, and students in each project and paper. I work to assure that the co-author and co-investigator groups in which I work are diverse, and are inclusive of those at early career stages. 

Advancement of diversity and inclusion represents a guiding value in my work. At Cornell I would seek out new opportunities to make the academy a more open and welcoming career path. I would be particularly excited to work on this in the context of the development of new programs and efforts in accordance with the early development of the Brooks School of Public Policy.  





\end{document} 
